\uuid{hS9F}
\exo7id{4487}
\titre{exo7 4487}
\auteur{quercia}
\organisation{exo7}
\datecreate{2010-03-14}
\isIndication{false}
\isCorrection{false}
\chapitre{Série numérique}
\sousChapitre{Familles sommables}
\cpge{mpsi-mp2i.raisonnement-vocabulaire-ensembliste.b, pcsi.raisonnement-vocabulaire-ensembliste.b, ptsi.raisonnement-vocabulaire-ensembliste.b, mpsi-mp2i.raisonnement-vocabulaire-ensembliste.c, pcsi.raisonnement-vocabulaire-ensembliste.d, ptsi.raisonnement-vocabulaire-ensembliste.d, mpsi-mp2i.denombrement.a, pcsi.denombrement.a, ptsi.denombrement.a, mp-mpi.variables-aleatoires-discretes.a, pc.familles-sommables, psi.denombrables-familles-sommables, pt.probabilites-lois-usuelles.a}
\module{Analyse}
\niveau{L1}
\difficulte{}

\contenu{
\texte{
$A$ étant un ensemble infini dénombrable, les ensembles suivants sont-ils
dénombrables~:
}
\begin{enumerate}
    \item \question{${\cal P}(A)$~?}
    \item \question{$\{\text{parties finies de }A\}$~?}
    \item \question{$\{\text{suites périodiques à valeurs dans }A\}$~?}
    \item \question{$\{\text{suites ultimement périodiques à valeurs dans }A\}$~?}
    \item \question{$\{\text{relations d'ordre total sur } A\}$~?}
\end{enumerate}
}
